\documentclass[11pt]{article}
\usepackage[a4paper,margin=0.9in]{geometry}
\usepackage[colorlinks=true,linkcolor=blue]{hyperref}
\usepackage{amsmath,amssymb}
\usepackage{booktabs}
\usepackage{caption}
\usepackage{graphicx}
\usepackage{longtable}
\usepackage{titlesec}
\usepackage{xcolor}

\title{Array Geometry Comparison: Theoretical Screening Results\\
\large IYPT Pinhole Sunglasses Project --- Array Design Sub-Study}
\author{}
\date{\today}

\begin{document}
\maketitle

\section{Purpose and Scope}
This report consolidates the theoretical (closed-form) screening comparison of
nine pinhole-array geometry families: square, hexagonal, rectangular (anisotropic),
concentric rings, ring-offset concentric, Fibonacci spiral, golden-angle
phyllotaxis, Archimedean spiral, and two disordered families (Poisson-disk,
jittered lattice).

At this stage, hole diameter $a_{\rm hole}$ is \emph{fixed} per myopia level $M$
from the project's existing acuity-optimal aperture chart
(\texttt{new\_sweeps/per\_M\_optimum\_map.csv}) and is not re-optimized here. This
study evaluates only geometry (pattern, hole count $N$, spacing) against two
independently validated closed-form physical constraints: ghost-image separation
and gaze-angle vignetting. No brightness/transmission objective is used in this
pass (explicit scope decision --- see Section~5).

\section{Data Integrity}
A degeneracy defect was found and corrected during this study: severe
overlap-rejection collapse (requested pitch smaller than hole diameter) produced
candidate rows where the per-family area formula no longer represented a valid
periodic tiling or convex hull, yielding physically impossible packing-density
values (open-area fraction $>1$). This was diagnosed to two independent
mechanisms --- a spacing-independent unit-cell infeasibility for lattice families,
and a boundary hole-radius omission for hull/circle-based families --- and fixed
at the root of the dataset, not patched downstream. The exclusion logic
(\texttt{degeneracy\_filter.py}) is now the single shared source of truth used
by every downstream script.

\input{report_tables/table1_integrity.tex}

\vspace{0.3cm}
\noindent Exclusion reasons (220 rows total, grouped):
\input{report_tables/table2_exclusion_reasons.tex}

\section{Safety Gates}
Two independent, composable gates are applied to every candidate row:
\begin{itemize}
  \item \textbf{Gate 1 (ghost safety):} worst-case ghost-separation margin $>0$,
  using the empirically validated formula $\Delta r = sM/(P_0-M)$
  ($R^2=0.995$ against direct two-hole coherent wave-optics propagation).
  \item \textbf{Gate 2 (vignetting safety):} worst-case margin $>0$ across
  gaze angle $\varphi \in [-25^\circ, 25^\circ]$, using the validated
  Handbook~2 Ch.~13 circle-overlap geometry.
\end{itemize}
A third gate (seed-reproducibility, for the two stochastic families) was
considered and explicitly \emph{not} applied in this pass, by decision: a
stochastic family's greater susceptibility to failing Gates~1--2 is intended to
show up naturally in survivor counts rather than be pre-filtered.

\subsection{Survivor Funnel per Family}
\input{report_tables/table3_funnel.tex}

\begin{figure}[h]
\centering
\includegraphics[width=0.95\textwidth]{report_figs/fig_funnel_survivors.png}
\caption{Rows surviving Gate 1 alone versus both gates, per family, across the
full cleaned dataset (6395 rows, 9 families).}
\end{figure}

\subsection{Ghost/Vignetting Divergence}
A central finding of this study: passing Gate~1 with an infinite margin (i.e.\
no hole pairs simultaneously active under the pupil at $\varphi=0$) does
\emph{not} imply Gate~2 is also satisfied. The two criteria must be evaluated
independently for every family --- one cannot be used as a proxy for the other.

\input{report_tables/table4_divergence.tex}

\begin{figure}[h]
\centering
\includegraphics[width=0.85\textwidth]{report_figs/fig_divergence.png}
\caption{Fraction of $\varphi{=}0$ ghost-safe candidates that fail vignetting
safety once gaze angle is integrated, per family. Divergence ranges from
37.8\% (concentric) to 88.9\% (Fibonacci spiral / golden angle).}
\end{figure}

\section{Comparison Metrics}
Two metrics are computed over the Gate~1$\cap$Gate~2 survivor set:

\begin{description}
  \item[Metric 1 --- maximum packing density.] Per $M$ (and a global,
  worst-case-across-$M$ scope), the single surviving candidate with highest
  $N_{\rm actual}$ per unit packed area.
  \item[Metric 2 --- Pareto front (density vs.\ binding margin).] The full
  2D dominance front over (packing density, $\min$(ghost margin, vignetting
  margin)), showing the trade-off rather than collapsing to one point.
\end{description}

\subsection{Metric 1 Winners (post-cleanup, all $N_{\rm actual}$ verified
non-degenerate)}
\input{report_tables/table5_metric1_winners.tex}

\subsection{Metric 2}
At the current tested spacing resolution, Metric~2's Pareto front collapses to
a single point at every $M$ (verified genuine via a synthetic-data control
that confirmed the dominance logic itself produces multi-point fronts when the
underlying data has a live density/margin trade-off). This reflects that
safety in the present dataset is largely binary --- trivially-safe (zero
active pairs) versus clearly unsafe --- rather than a graded trade-off region.
Widening the tested spacing range toward the safety boundary is the natural
next probe if a live Pareto trade-off is worth resolving explicitly.

\section{Honest Scope Notes}
\begin{itemize}
  \item Hole diameter is fixed from the acuity-optimal chart, not searched ---
  this study answers ``which geometry, at the size already known to be
  acuity-optimal,'' not ``which size and geometry jointly.''
  \item No brightness/transmission constraint is applied in this pass, by
  explicit decision, so that reproducibility/manufacturability
  considerations for the two stochastic families could be observed
  without being pre-filtered by an unrelated gate.
  \item This study uses the validated \emph{closed-form} ghost and
  vignetting formulas only --- not the full multi-hole wave-optics (ASM)
  engine. A subsequent confirmation pass on the Metric~1/Metric~2 finalists
  using the full validated propagation engine is the recommended next step
  before any physical fabrication decision.
  \item \texttt{family\_ranking.csv} (an earlier, independent output) was
  checked and confirmed to draw only on ghost/vignetting-margin columns, not
  density --- it was not affected by the degeneracy defect described in
  Section~2.
\end{itemize}

\input{confirmation_addendum_section.tex}

\section{Recommended Next Step}
With all 9 Metric-1 finalists confirmed under full-physics propagation
(Section~7), the remaining gap is between simulation and reality: no design
in this study has yet been checked against a physically fabricated sample.
Immediate next steps: (1) regenerate \texttt{new\_sweeps/blur\_sweep.csv} at
$N{=}4096$ project-wide, not just for the 9 finalists spot-checked here; (2)
select 2--3 finalists spanning the tested $M$ range for laser-cut shim
fabrication, prioritizing manufacturability margin over raw blur/density
optimality; (3) measure blur, transmission, and gaze-angle ghosting on the
fabricated prototypes using the project's established phone-camera afocal
photography protocol, and compare directly against this report's and
Section~7's predictions.

\end{document}
